\begin{solapartat}
\figura[0.45]{sol_fig1.png}

El moment angular de Mercuri respecte el centre del Sol s'expressa com:
\[ \vec{L}_{Mer,Sol} = m\vec{r}\times\vec{v}, \]
on $\vec{r}$ és el vector de posició que apunta del centre del Sol a Mercuri, $m$ és la massa de Mercuri i $\vec{v}$ és la seva velocitat. El mòdul del moment angular de Mercuri respecte el Sol, s'expressa com:
\[ L_{Mer,Sol} = m\,r\,v\sin\theta, \]
on $\theta$ és l'angle que formen $\vec{r}$ i $\vec{v}$. Com al periheli i a l'afeli $\vec{r}$ i $\vec{v}$ són perpendiculars, $\theta = \pi/2$ i $\sin\theta = 1$. Per tant, si imposem la conservació del moment angular entre el periheli i l'afeli tenim:
\[ m\,r_a\,v_a = m\,r_p\,v_p \]
I si aïllem el mòdul de la velocitat de Mercuri al periheli obtenim:
\[ v_p = v_a\frac{r_a}{r_p} = 3,88\cdot 10^{4}\ \mathrm{m/s}\,\frac{6,99\cdot 10^{13}\ \mathrm{m}}{4,60\cdot 10^{13}\ \mathrm{m}} = 5,90\cdot 10^{4}\ \mathrm{m/s} \]
Com el moment angular és constant, $L_{Mer,Sol} = m\,r\,v\sin\theta = \mathit{constant}$, i tant la velocitat com l'angle varien, llavors el mòdul de la velocitat no és constant.

Com el moment angular i la massa són constants:
\[ v = \frac{\mathit{constant}}{r\sin\theta} \]
La velocitat serà mínima en el punt on el producte $r\sin\theta$ sigui màxim, i aquest punt és l'afeli atès que és el punt de l'òrbita s'allunya més del Sol i el sinus de l'angle pren el seu valor màxim $\sin\theta = 1$.
\end{solapartat}
